% Generated draft: requires mathematical and source review.
\documentclass[11pt]{article}
\usepackage[margin=25mm]{geometry}
\usepackage{fontspec}
\setmainfont{Latin Modern Roman}
\usepackage{amsmath,amssymb,mathtools,unicode-math}
\setmathfont{Latin Modern Math}
\usepackage{enumitem,needspace,xurl,hyperref}
\setlength{\parindent}{0pt}
\setlength{\parskip}{5pt}
\setlength{\emergencystretch}{3em}
\newcommand{\N}{\mathbb{N}}
\newcommand{\Z}{\mathbb{Z}}
\newcommand{\Q}{\mathbb{Q}}
\newcommand{\R}{\mathbb{R}}
\newcommand{\C}{\mathbb{C}}
\newcommand{\F}{\mathbb{F}}
\newcommand{\A}{\mathbb{A}}
\newcommand{\PP}{\mathbb{P}}
\newcommand{\E}{\mathbb{E}}
\newcommand{\eps}{\varepsilon}
\newcommand{\dd}{\,\mathrm{d}}
\begin{document}
\section*{UnsolvedMath ID: 30004143}
\textbf{Source problem code:} catalogue entry
\subsection{Definitions and mathematical statement}
Write \(\Gamma=(G,\ell,w)\) for a connected abstract tropical curve represented by a weighted finite metric graph. Here \(G\) is a finite connected graph, \(\ell\) records its edge lengths, and \(w\) denotes its weight data. In the usual finite-edge metric convention, \(\ell(e)\in\R_{>0}\) for each edge \(e\). The supplied record does not specify the meaning or domain of the weights, or whether unbounded legs or marked points are permitted; these conventions must be fixed before the requested theory can be formulated precisely.

Let \(|G|\) be the topological realization of the underlying finite graph and choose \(v\in |G|\). Its ordinary topological fundamental group is
\[
\pi_1(|G|,v)\cong F_{b_1(G)},\qquad b_1(G)=|E(G)|-|V(G)|+1,
\]
where \(F_r\) is the free group of rank \(r\). This construction concerns the underlying topological graph and does not by itself specify how tropical weights, ramification, or algebraic realization data should enter the requested theory.

A graph of groups on \(G\) consists of a group \(A_x\) at each vertex \(x\), a group \(A_e\) at each unoriented edge \(e\), and an injective homomorphism from \(A_e\) to the group at each incident endpoint, with both incidences included for a loop. Once these data are chosen, one can form the associated graph-of-groups fundamental group. The source does not prescribe these groups or homomorphisms, or how they should be obtained from \(\Gamma\).

\textbf{Requested research construction.} For each connected abstract tropical curve represented by a weighted finite metric graph, formulate an appropriate fundamental-group theory, including graph-of-groups constructions, with classification statements for the intended covers, harmonic morphisms, and realizations of admissible covers. A precise formulation must identify the objects and morphisms to be classified, the additional data attached to \(\Gamma\), and the sense in which the proposed group or other fundamental-group object classifies them.

\textbf{Unresolved qualifications.} The original statement lists covers, harmonic morphisms, \emph{or} realizations of admissible covers, whereas the supplied clean statement requires covers, harmonic morphisms, \emph{and} realizations of admissible covers. It is therefore unclear whether these are alternative targets or simultaneous requirements. The record supplies no definitions of the relevant covering categories, permitted ramification, harmonicity conditions, admissibility conditions, or algebraic realization setting. Consequently, this is a research-program statement with unresolved specifications, rather than a fully determined conjecture asserting the existence of a particular group or classification equivalence.
\subsection{Short English statement}
Seek a fundamental-group theory for connected weighted finite metric graphs viewed as tropical curves, including graph-of-groups constructions that classify suitable covers and related morphisms. The ordinary graph fundamental group is available, but the source leaves the extra tropical data and intended classification categories unspecified. The original and cleaned statements also differ on whether the listed classification targets are alternatives or simultaneous requirements. This statement-only draft requires human mathematical review.
\subsection{Sources}
\begin{itemize}
\item Supplied local UnsolvedMath registry record, permanent ID \(30004143\), title \emph{Fundamental Groups of Tropical Curves}, source problem code \texttt{catalogue entry}, in the Oberwolfach Reports Open Problems collection. No individual proposer is supplied; the recorded proposed year is 2019.
\item Supplied report citation: \emph{Tropical Geometry: new directions} (2020), Oberwolfach Reports, pp.\ 1245--1308, DOI \texttt{10.4171/owr/2019/21}. \url{https://doi.org/10.4171/owr/2019/21}. The supplied background describes Oberwolfach Reports as documenting mathematical workshops and their open-problem discussions. No problem-specific page or individual attribution is supplied.
\item Supplied literature source, labelled ``Tropical geometry problem report'': \url{https://ems.press/content/serial-article-files/46800}.
\item Supplied literature source, labelled ``Non-Archimedean graph-of-groups fundamental groups'': \url{https://link.springer.com/article/10.1007/s00029-021-00651-4}.
\item Literature assessment date: \texttt{2026-08-22}. This is supplied dated curation, not an independently verified assessment. The supplied status is \emph{open}. The supplied assessment states that the underlying graph has an ordinary free fundamental group, that graph-of-groups, stack, and orbifold interpretations can yield different fundamental groups, and that a unique general notion answering the question has not been established. The background explicitly requests independent verification of time-sensitive status and cited claims.
\item The registry labels the statement \texttt{corrected\_verified} and reports a manual semantic, typographic, and LaTeX review confirming that the clean statement preserves the intended target. These are supplied provenance claims; they do not resolve the formulation issues recorded here or constitute independent mathematical verification.
\end{itemize}
\end{document}
